\documentclass[10pt,a4paper]{amsart}
\usepackage[margin=19mm]{geometry}
\usepackage{amsmath,amssymb,mathtools}
\usepackage[scr=rsfs,cal=euler]{mathalfa}
\usepackage{xfrac}
\usepackage[unicode,hidelinks,bookmarks=false]{hyperref}
\newcommand{\ie}{i.e.\ }
\newcommand{\cf}{cf.\ }
\newcommand{\resp}{respectively\ }
\newcommand{\Subsection}{\subsection}
\providecommand{\nobreakdash}{\nobreak\hbox{-}}
\setlength{\emergencystretch}{2em}
\multlinegap0pt
\usepackage{bm}
\usepackage{bbm}
\usepackage{dsfont}
\usepackage{xspace}
\usepackage{cancel}
\usepackage{mathtools}
\usepackage{amsthm}
\makeatletter
\@ifundefined{lemma}{\newtheorem{lemma}{Lemma}}{}
\makeatother
\newcommand{\malA}{\mathcal{A}}
\newcommand{\malE}{\mathcal{E}}
\newcommand{\mbfH}{\mathbf{H}}
\newcommand{\mbfL}{\mathbf{L}}
\newcommand{\p}{\partial}
\title[Error estimates of linear decoupled structure-preserving inc]{Error estimates of linear decoupled structure-preserving incremental viscosity splitting methods for the Cahn--Hilliard--Navier--Stokes system}
\author{Baolin Kuang, Hongfei Fu, Xiaoli Li}
\date{Source: arXiv:2508.01141v1 (CC BY 4.0)}
\begin{document}
\maketitle

\noindent\emph{Source and licence.} Error estimates of linear decoupled structure-preserving incremental viscosity splitting methods for the Cahn--Hilliard--Navier--Stokes system. Version arXiv:2508.01141v1 (CC BY 4.0), distributed under the Creative Commons Attribution 4.0 International licence, as recorded in the arXiv licence metadata for this exact version and shown on the abstract page https://arxiv.org/abs/2508.01141v1. Full rights notice: licenses/arxiv-2508.01141-CC-BY-4.0.txt.\par\smallskip

\noindent\textit{Editorial scope.} Counted unit: the Lemma whose source label is \texttt{lem:err:CH} in the article (source0.tex lines 691--702, source label \texttt{lem:err:CH}), delivered together with the article's own complete original proof of it (source0.tex lines 704--834, one \texttt{proof} environment of 131 lines). No mathematics is written, repaired or re-derived: statement and proof are the article's own text line for line. Required context retained uncounted: model:sav:mu (lines 245--254); model:sav:phi (lines 245--254); scheme:1th:mu (lines 274--281); scheme:1th:phi (lines 274--281); bound:solu (lines 482--484); assump:e1 (lines 633--638); bound\_phi (lines 656--658). Every definition, display and label of those lines is retained unchanged. Omitted from this package and not redistributed: 1--244, 255--273, 282--481, 485--632, 639--655, 659--690, 703--703, 835--1843. Nothing inside a retained excerpt is omitted or altered: every retained body is delivered exactly as the article's lines are written, so no omission receipt of any kind accompanies this package. Citations: the retained excerpts carry no citation command at all, so no bibliography entry of the article is transcribed and no reference is redistributed. Reference adaptation: none was needed and none was made; every reference inside the delivered unit resolves inside it, so no unresolved reference or citation is produced. Printed slips kept literal: no printed word, symbol, hypothesis or number of the counted statement or of its proof was altered. Layout and class: the article's own class is \texttt{svjour3} with packages including fix-cm, mathptmx, amsmath, amssymb, amsfonts, graphicx, subfigure, stmaryrd, fontenc, latexsym, mathrsfs, float, xcolor, algorithm; the wrapper is the lane's pinned \texttt{amsart} editorial wrapper at 10pt on A4, which restates the theorem-like environments the retained text needs and adds the article's own macro definitions that the retained text needs (\texttt{\textbackslash malA}, \texttt{\textbackslash malE}, \texttt{\textbackslash mbfH}, \texttt{\textbackslash mbfL}, \texttt{\textbackslash p}), with extra wrapper packages (bm, bbm, dsfont, xspace, cancel, mathtools). The delivered numbering is the unit's own. Assets: no retained span contains a figure environment, a graphics-inclusion command, a PDF-page inclusion, a table or a TikZ picture; no image file is copied into this package, the package declares no asset dependency and no image file is redistributed with it. Licence: version arXiv:2508.01141v1 of this article is distributed under the Creative Commons Attribution 4.0 International licence, as recorded in the arXiv licence metadata for this exact version and shown on the abstract page https://arxiv.org/abs/2508.01141v1. A full rights notice is delivered at licenses/arxiv-2508.01141-CC-BY-4.0.txt. Characters: every retained excerpt is pure ASCII, so the delivered engine, XeLaTeX with its default Latin Modern text font, reports no missing character.
\begin{align}
	& \phi_t+ \xi(t) \nabla \cdot (\bm u \phi) =M \Delta \mu,  \label{model:sav:phi} \\
	& \mu=- \lambda\Delta \phi +\lambda  \beta \phi+ \lambda \xi(t) F^{\prime}(\phi),  \label{model:sav:mu}\\
	& \bm u_t +  \xi(t) \bm u\cdot\nabla \bm u-\nu \Delta \bm u +\nabla p=-  \xi(t) (\phi \nabla \mu + \gamma(t) \nabla p) + \gamma(t)  \nabla p,  \label{model:sav:u}\\
	& \nabla \cdot \bm u = 0,    \label{model:sav:divu} \\
	&  R_t =\frac{1}{2  \sqrt{E_1(\phi)+\delta_0}}   ( F^{\prime}(\phi), \phi_t)
	+  \frac{1}{2 \lambda  \sqrt{E_1(\phi)+\delta_0}} \big((\mu , \nabla \cdot (\bm u \phi) )+  ( {\bm u} ,  \phi  \nabla\mu ) \big) \notag\\
	&\qquad  + \frac{1}{2 \lambda  \sqrt{E_1(\phi)+\delta_0}} \bm b(\bm u, \bm u, \bm u) + \frac{ \gamma(t)}{2 \lambda  \sqrt{E_1(\phi)+\delta_0}}  (\nabla p, \bm u), 
	\label{model:sav:r}
\end{align}

\begin{align}
	& \p_{\tau} \phi^{n+1} + \xi^{n+1}\nabla \cdot({\bm u}^n \phi^n  )=M \Delta \mu^{n+1}, \quad  
	\p_{\bm n} \phi^{n+1}|_{\partial \Omega}=0, \label{scheme:1th:phi}\\
	& \mu^{n+1}=- \lambda \Delta \phi^{n+1} + \lambda \beta \phi^{n+1} + \lambda  \xi^{n+1} F^{\prime}(\phi^n), \quad   
	\p_{\bm n} \mu^{n+1} |_{\partial \Omega}=0, \label{scheme:1th:mu} \\
	& \frac{\widehat {\bm u}^{n+1} -  \bm u^n}{\tau} + \xi^{n+1} {\bm u}^n\cdot\nabla {\bm u}^n - \nu\Delta   \widehat {\bm u}^{n+1} = -\xi^{n+1}  \big( \phi^n \nabla \mu^n+\gamma^n \nabla p^n \big), ~
	\widehat{\bm u}^{n+1}|_{\partial \Omega}=0, \label{scheme:1th:uhat}   
\end{align}

	\begin{equation} \label{bound:solu}
		\|\bm u^{n+1}\|^2 +\|\phi^{n+1}\|_1^2 +|R^{n+1}|^2  +   {\nu \tau\sum_{k=0}^n\|\nabla {\bm u}^{k+1} \|^2}\leq K_1,
	\end{equation}

\begin{equation} \label{assump:e1}
	\begin{aligned}
		& \phi \in H^{2}(0,T; L^2(\Omega)) \cap W^{1, \infty}(0,T; H^1(\Omega)) \cap L^{\infty}(0,T; H^2(\Omega)), ~  \mu \in W^{1, \infty}(0,T; H^2(\Omega)),\\
		& \bm u \in H^{2}(0,T; \mbfL^2(\Omega)) \cap W^{1, \infty}(0,T;\mbfH^2(\Omega)\cap   \mbfH_0^1(\Omega)), ~ p \in W^{1, \infty}(0,T; L^2(\Omega))\cap L^{\infty}(0,T;H^1(\Omega)).
	\end{aligned}  
\end{equation}

\begin{equation} \label{bound_phi}
	\|\phi^n\|_{L^\infty} \le C_{\phi}, \quad 0 \leq n \le N_t.
\end{equation}

\begin{lemma} \label{lem:err:CH}
	Under the assumption \eqref{assump:e1}, we have
	\begin{equation*}  
		\begin{aligned}
			&\frac{\lambda }{2}\p_\tau \|\nabla e_\phi^{n+1}\|^2+\frac{1+\beta}{2}\p_\tau \|e_\phi^{n+1}\|^2
			+\frac{3M}{4}\|e_\mu^{n+1}\|_1^2 \\
			& \leq  -\lambda  \frac{e_R^{n+1}}{\sqrt{E_1(\phi^{n})+\delta_0}} \big(F^{\prime}(\phi^{n}),\p_\tau  e_\phi^{n+1}\big)
			+C|e_R^{n+1}|^2 + C\|e_\phi^{n+1}\|_1^2 + C\|e_\phi^n\|_1^2			\\
			& \quad    +C\|e_{\bm u}^n\|^2+C \tau \big(\|\bm u_t\|_{L^2(J_n ; \mbfL^2(\Omega))}^2+\|\phi_t\|_{L^2(J_n ; H^1(\Omega))}^2+\|\phi_{tt}\|_{L^2(J_n ; L^2(\Omega))}^2\big).
		\end{aligned}
	\end{equation*}
\end{lemma}

\begin{proof} First, subtracting \eqref{scheme:1th:phi} from \eqref{model:sav:phi} at $t=t_{n+1}$, we obtain the error equation
	\begin{equation} \label{err:ephi}
		\p_\tau  e_\phi^{n+1}-M \Delta e_\mu^{n+1}=R_\phi^{n+1}+E_N^{n+1},
	\end{equation}
	where  
	\begin{equation*} 
		\begin{aligned}
			& R_\phi^{n+1}:=\p_\tau \phi(t_{n+1})- \phi_t(t_{n+1})=\frac{1}{\tau} \int_{t_n}^{t_{n+1}}(t-t_n) \phi_{tt} d t, \\
			& E_N^{n+1}:=\frac{R^{n+1}}{\sqrt{E_1(\phi^n)+\delta_0}} \nabla \cdot({\bm u}^n \phi^n) -\frac{R(t_{n+1})}{\sqrt{E_1(\phi(t_{n+1}))+\delta_0}}\nabla \cdot  (\bm u(t_{n+1})\phi(t_{n+1})).
		\end{aligned} 
	\end{equation*}
	Taking the inner products of \eqref{err:ephi} with $e_\mu^{n+1}$ and $e_\phi^{n+1}$, respectively, we have
	\begin{equation} \label{err:ephi:1}
		\big(\p_\tau  e_\phi^{n+1}, e_\mu^{n+1}\big)+M\|\nabla e_\mu^{n+1}\|^2= (R_\phi^{n+1}, e_\mu^{n+1})+(E_N^{n+1}, e_\mu^{n+1}),
	\end{equation}
	\begin{equation} \label{err:ephi:2}
		\frac{1}{2}\p_\tau \| e_\phi^{n+1}\|^2+ \frac{\tau}{2}\|\p_\tau  e_\phi^{n+1}\|^2=(R_\phi^{n+1}, e_\phi^{n+1})+(E_N^{n+1}, e_\phi^{n+1})-M(\nabla e_\mu^{n+1}, \nabla e_\phi^{n+1}).
	\end{equation}
	
	Next, subtracting \eqref{scheme:1th:mu} from \eqref{model:sav:mu} at $t=t_{n+1}$, we get
	\begin{equation}  \label{err:emu}
		e_\mu^{n+1}=- \lambda\Delta e_\phi^{n+1}+ \lambda \beta e_\phi^{n+1} + \lambda \frac{e_R^{n+1}}{\sqrt{E_1(\phi^{n}))+\delta_0}} F^{\prime}(\phi^{n})+E_F^{n+1},
	\end{equation} 
	where  
	\begin{equation*}  
		E_F^{n+1}:=\lambda  \Big(\frac{F^{\prime}(\phi(t_{n+1}))}{\sqrt{E_1(\phi(t_{n+1}))+\delta_0}}-\frac{F^{\prime}(\phi^n)}{\sqrt{E_1(\phi^n)+\delta_0}}\Big) R(t_{n+1}).
	\end{equation*} 
	Taking the inner products of \eqref{err:emu} with $M e_\mu^{n+1}$ and $\p_\tau e_\phi^{n+1}$, respectively, we obtain
	\begin{equation} \label{err:emu:1}
		\begin{aligned}
			M\|e_\mu^{n+1}\|^2 
			& =    M \lambda(\nabla e_\mu^{n+1}, \nabla e_\phi^{n+1}) +M  \lambda 
			\beta(e_\phi^{n+1},e_\mu^{n+1}) \\
			&\qquad +M \lambda  \frac{ e_R^{n+1}}{\sqrt{E_1(\phi^{n})+\delta_0}}(F^{\prime}(\phi^{n}), e_\mu^{n+1})+M(E_F^{n+1}, e_\mu^{n+1}),
		\end{aligned}
	\end{equation}
	\begin{equation} \label{err:emu:2}
		\begin{aligned}
			\big( e_\mu^{n+1},\p_\tau  e_\phi^{n+1}\big)
			& =  \frac{\lambda }{2}\big( \p_\tau \|\nabla e_\phi^{n+1}\|^2 + \tau\|\p_\tau  \nabla e_\phi^{n+1}\|^2\big)  
			+  \frac{\lambda \beta}{2}\big( \p_\tau\| e_\phi^{n+1}\|^2 + \tau\|\p_\tau   e_\phi^{n+1}\|^2\big)\\
			& \qquad  +\big(E_F^{n+1}, \p_\tau  e_\phi^{n+1}\big) + \lambda \frac{e_R^{n+1}}{\sqrt{E_1(\phi^{n})+\delta_0}} \big(F^{\prime}(\phi^{n}),\p_\tau  e_\phi^{n+1}\big)   .
		\end{aligned}
	\end{equation}
	
	
	Now, adding \eqref{err:ephi:1}--\eqref{err:ephi:2} and \eqref{err:emu:1}--\eqref{err:emu:2} together, we obtain
	\begin{equation} \label{err:CH:all:e1}
		\begin{aligned}
			&\frac{\lambda }{2}  \p_\tau \|\nabla e_\phi^{n+1}\|^2 %+ \tau\|\p_\tau  \nabla e_\phi^{n+1}\|^2\big)
			+\frac{\lambda (1+\beta)}{2}  \p_\tau\| e_\phi^{n+1}\|^2 %+ \tau\|\p_\tau   e_\phi^{n+1}\|^2\big)   
			+ M\| e_\mu^{n+1}\|_1^2 \\
			&  \le - \lambda \frac{e_R^{n+1}}{\sqrt{E_1(\phi^{n})+\delta_0}} \big(F^{\prime}(\phi^{n}),\p_\tau  e_\phi^{n+1}\big)+\lambda M   \frac{e_R^{n+1}}{\sqrt{E_1(\phi^{n})+\delta_0}}(F^{\prime}(\phi^{n}), e_\mu^{n+1}) \\
			& \quad +\big(E_F^{n+1}, M e_\mu^{n+1}-\p_\tau  e_\phi^{n+1}\big)+ (E_N^{n+1}, e_\mu^{n+1}+e_\phi^{n+1}) \\
			& \quad  +(R_\phi^{n+1}, e_\mu^{n+1})+(R_\phi^{n+1}, e_\phi^{n+1})+ \lambda M\beta(e_\phi^{n+1},e_\mu^{n+1})
			=:\sum\limits_{i=1}^7\malA_i.
		\end{aligned}
	\end{equation}
	
	We pay attention to the estimates of the right-hand side of \eqref{err:CH:all:e1}. First, the second term can be bounded by
	\begin{equation*} 
		\malA_2\leq   \frac{M}{16}\|e_\mu^{n+1}\|^2+C|e_R^{n+1}|^2,
	\end{equation*}
	where the uniform boundedness of $\phi^n$ is utilized. Second, the last three terms can be estimated in a standard way that 
	\begin{equation*} 
		\malA_{5} + \malA_{6} +  \malA_{7}
		\le \frac{M}{16}\|e_\mu^{n+1}\|^2+ C\| e_\phi^{n+1}\|^2+ C\tau \|\phi_{tt}\|_{L^2(J_n ; L^2(\Omega))}^2,
	\end{equation*}
	where	$\|R_\phi^{n+1}\|^2 \le \tau \|\phi_{tt}\|_{L^2(J_n ; L^2(\Omega))}^2$ 	is applied.
	
	Third, owing to \eqref{err:ephi}, the third term can be rewritten as
	\begin{equation*}
		\begin{aligned}
			\malA_3
			&=\big(E_F^{n+1}, M e_\mu^{n+1}-M \Delta e_\mu^{n+1}-R_\phi^{n+1}-E_N^{n+1}\big) \\
			& =\big(E_F^{n+1}, M e_\mu^{n+1}-R_\phi^{n+1}-E_N^{n+1}\big)+M\big(\nabla E_F^{n+1}, \nabla e_\mu^{n+1}\big). %  =: \sum\limits_{j=1}^2\malA_{3,j}.
		\end{aligned}
	\end{equation*}
	Note that $E_N^{n+1}$ is equivalent to  
	\begin{equation} \label{err:EN:all}
		\begin{aligned}
			E_N^{n+1} 
			& =  -\frac{R(t_{n+1})}{\sqrt{E_1(\phi(t_{n+1}))+\delta_0}}  
			\nabla \cdot\Big( \phi(t_n)\int_{t_n}^{t_{n+1}} \bm u_t dt+\bm u (t_{n+1}) \int_{t_n}^{t_{n+1}} \phi_t dt \Big) \\
			& \qquad  + \malE_R^{n+1} \nabla \cdot{\left(\bm u (t_{n}) \phi(t_{n})\right)} -\frac{R^{n+1}}{\sqrt{E_1(\phi^n)+\delta_0}} \nabla \cdot\left(\bm u (t_{n}) e_\phi^n+e_{\bm u }^n \phi^n\right),
		\end{aligned}
	\end{equation}
	where 
	\begin{equation} \label{err:Er}
		\malE_R^{n+1}:=\frac{R^{n+1}}{\sqrt{E_1(\phi^n)+\delta_0}}-\frac{R(t_{n+1})}{\sqrt{E_1(\phi(t_{n+1}))+\delta_0}} 
		\le C|e_R^{n+1}|+ C|e_\phi^n|  +C\int_{t_n}^{t_{n+1}}|\phi_t| dt.
	\end{equation}
	Besides, the boundedness of $\phi(t_{n+1})$ and $\phi^n$ imply that
	\begin{equation}  \label{err:E_F}
		\|E_F^{n+1}\|_1^2 \le		C\|e_\phi^n\|_1^2  +C \tau \|\phi_t\|^2_{L^2(J_n ; H^1(\Omega))}.
	\end{equation}
	Therefore, using the integration by parts formula and \eqref{err:EN:all}--\eqref{err:E_F}, we see 
	\begin{equation} \label{err:EN:EF}
		\begin{aligned}
			 (E_F^{n+1}, E_N^{n+1}) 	& = \frac{R(t_{n+1})}{\sqrt{E_1(\phi(t_{n+1}))+\delta_0}}
			\big(\nabla E_F^{n+1},  \phi(t_n)\int_{t_n}^{t_{n+1}} \bm u_t dt+\bm u (t_{n+1}) \int_{t_n}^{t_{n+1}} \phi_t dt\big) \\
			& \quad - \malE_R^{n+1} \big(\nabla E_F^{n+1},  \bm u(t_{n}) \phi(t_{n})\big) 
			+\frac{R^{n+1}}{\sqrt{E_1(\phi^n)+\delta_0}}\big(\nabla E_F^{n+1}, \bm u(t_{n}) e_\phi^n +e_{\bm u}^n \phi^n\big)\\
			&  \leq C|e_R^{n+1}|^2 +C\|e_\phi^n\|_1^2+C\|e_{\bm u}^n\|^2
			+C\tau\big( \|\phi_t\|_{L^2(J_n ; H^1(\Omega))}^2 + \|\bm u_t\|_{L^2(J_n; \mbfL^2(\Omega))}^2 \big),
		\end{aligned}
	\end{equation}
	where the boundedness of \eqref{bound:solu}  and \eqref{bound_phi} are also applied in the last step.
	
	Thus, combing \eqref{err:E_F} and \eqref{err:EN:EF} together and applying the Cauchy-Schwarz inequality, it yields
	\begin{equation*} 
		\begin{aligned}
			\malA_{3} 
			&  \leq C|e_R^{n+1}|^2+\frac{M}{8}\|e_\mu^{n+1}\|_1^2
			+C\|e_{\bm u}^n\|^2+C\|e_\phi^n\|_1^2\\ %+ C\|\nabla e_\phi^n\|^2\\
			& \quad +C \tau \big(\|\phi_t\|_{L^2(J_n; H^1(\Omega))}^2
			+\|\phi_{tt}\|_{L^2(J_n; L^2(\Omega))}^2+\|\bm u_t\|_{L^2(J_n; \mbfL^2(\Omega))}^2\big).
		\end{aligned}
	\end{equation*}
	
	Finally, by using integration by parts formula again and similar to the estimate of \eqref{err:EN:EF}, we have
	\begin{equation*} 
		\begin{aligned}
			\malA_{4} 
			& \leq C|e_R^{n+1}|^2+C\|\nabla e_\phi^{n+1}\|^2+ \frac{M}{8}\|\nabla e_\mu^{n+1}\|^2 +C\|e_\phi^n\|^2+C\|e_{\bm u}^n\|^2\\
			& \quad +C \tau \big(\|\phi_t\|_{L^2(J_n ; L^2(\Omega))}^2+\|\bm u_t\|_{L^2(J_n ; \mbfL^2(\Omega))}^2\big).
		\end{aligned}
	\end{equation*}
	
Now, inserting the estimates of $\malA_2-\malA_{7}$ into \eqref{err:CH:all:e1}, we obtain the desired result.
\end{proof}

\end{document}
